% !PRINTED 554
\weberpage{554}
%<<CONT
and $\chi$ an integral functional. We now choose (according to \S.\,143, 3.) a number $\beta$ divisible by $\chi$ in such a way that $\beta:\chi$ becomes relatively prime to $\omega$, and set

\weq{3}{\psi\chi = \beta.}

Since $\varphi:\psi$ is now a number, $\psi$ is equivalent to $\varphi$, and $\psi$ is a representative, coprime to $\omega$, of the class $C^1$).

\S.\,153.

\gesperrt{The class number of the field $\Omega$.}

We come now to the proof of the important theorem:

1. \gesperrt{The number of ideal classes of a field $\Omega$ is finite.}

This theorem is equivalent to the following:

2. \gesperrt{In every class there are integral functionals whose absolute norm does not exceed a certain finite number depending only on the nature of the field $\Omega$.}

For since every integral functional is a factor of its absolute norm, and every integer has only a finite number of non-associated factors, there are, apart from the associated ones, only a finite number of integral functionals that have a given number as absolute norm. If now it can be proved that in every class there occurs an integral functional whose absolute norm under a finite number

\begin{footnotesize}
$^1$) We wish here to point out in passing an analogy of the equivalence of ideals. The whole theory of algebraic numbers can be carried over, with the necessary modifications, to the theory of algebraic functions, which in turn finds its geometric expression in the theory of algebraic curves. To the ideals there then correspond systems of points on a fixed fundamental curve, and to the principal ideals full systems of intersection points of the fundamental curve with another algebraic curve. When two systems of points complete each other to a full system of intersection points, which would find its expression in the above formula $\varphi\chi = \alpha$, each of the two systems of points $\varphi$, $\chi$ is called the \gesperrt{residue} of the other. Two systems of points which, like $\varphi$, $\psi$, have the same residue are called corresidual in geometry. This concept thus corresponds to equivalence. (Cf. \gesperrt{Brill} u. \gesperrt{Nöther}, Mathem. Annalen, Bd. 7. \gesperrt{Salmon}, ``higher plane Curves'', deutsch von \gesperrt{Fiedler}.)
\end{footnotesize}
%CONT>>
